\section{Theory}
\label{sec:theory}

\subsection{Measurement model and Cram\'er--Rao bound}

Each test structure is a waveguide of bottom width $w$ in a silicon film of thickness $t$, read out by a ring
resonator and, in some of the sets we compare, also by a Mach--Zehnder interferometer (MZI). The unknown parameters form a vector $\mathbf{p}$ that always
contains the strip width offset, the film thickness and the strip sidewall angle $\theta$, measured from vertical
with the top narrower. A ring gives two observables, the effective index $\neff$ through the resonance position and
the group index $\ngr$ through the free spectral range; an MZI gives one phase. Reading $\neff$ from a resonance
position assumes that the resonance order and the ring length are known, so that a wavelength shift maps onto an index
change; we make this assumption throughout. Stacking all observables into
$\mathbf{y}=\mathbf{f}(\mathbf{p})+\boldsymbol{\epsilon}$ with Gaussian noise of covariance
$\boldsymbol{\Sigma}$, the Fisher information matrix is
\begin{equation}
\mathbf{F}=\mathbf{J}^{\mathsf T}\boldsymbol{\Sigma}^{-1}\mathbf{J}+\boldsymbol{\Lambda},\qquad
J_{kj}=\frac{\partial f_k}{\partial p_j},
\label{eq:fisher}
\end{equation}
where $\boldsymbol{\Sigma}=\mathrm{diag}(\sigma_k^2)$ in the baseline model (Section~\ref{sec:noise}
relaxes this) and $\boldsymbol{\Lambda}$ is diagonal in the parameters that
have a Gaussian prior, with entries $1/\sigma_{p,j}^2$ (a prior on a difference of two parameters adds the
corresponding rank-one term). The CRLB on the angle is~\cite{kay1993}
\begin{equation}
\sth\ge\sqrt{\left[\mathbf{F}^{-1}\right]_{\theta\theta}}.
\label{eq:crlb}
\end{equation}
When $\mathbf{F}$ is singular along a direction that has a component on $\theta$, no unbiased estimator has finite
variance and we call the angle unidentifiable ($\sth=\infty$). Numerically, $\mathbf{J}$ is column-normalized and
singular values below a fixed small fraction of the largest are treated as zero.

The noise is set by the precision $\sigma_\lambda$ with which a resonance wavelength can be located. Janz
\textit{et~al.}\ report that wavelength changes of \qty{1}{pm} or less are routinely measured with calibrated tunable
lasers, and that a temperature-stabilized ring reached a \qty{0.1}{pm} uncertainty set by the laser
repeatability~\cite{janz2023}. We take \qty{1}{pm} as the primary case and \qty{0.1}{pm} as the best case. The
conversion to the three observables is ours. A resonance shift obeys
$\delta\lambda/\lambda=\delta\neff/\ngr$, so $\sigma_{\neff}=\ngr\sigma_\lambda/\lambda$; the group index follows from
one free spectral range, $\ngr=\lambda^2/(L\,\mathrm{FSR})$, measured with two resonances, so
$\sigma_{\ngr}=\sqrt2\,\sigma_\lambda\ngr^2L/\lambda^2$; and an MZI fringe read as precisely as a ring resonance gives
$\sigma_\phi=2\pi\ngr L\sigma_\lambda/\lambda^2$. With $\lambda=\qty{1.55}{\micro m}$, a ring and MZI length
$L=\qty{100}{\micro m}$ and $\ngr=4.26$, the largest group index among the structures used, $\sigma_\lambda=\qty{1}{pm}$
gives $\sigma_{\neff}=\num{2.75e-6}$, $\sigma_{\ngr}=\num{1.07e-3}$ and $\sigma_\phi=\qty{1.11e-3}{rad}$. These
three values are derived, not measured. The MZI value assumes that a broad fringe is located as well as a narrow
resonance, and the baseline treats the three errors as independent. If $\neff$ is read from one of the two
resonances that define the FSR, the $\neff$ and $\ngr$ errors of a ring are correlated with $\rho=-1/\sqrt2$; with
the mean of the two resonances they are uncorrelated. Section~\ref{sec:noise} evaluates both cases, a common
calibration offset and temperature drift. At \qty{0.1}{pm} the ring must also be held to about \qty{1.25}{mK}, since its resonance moves
by roughly \qty{80}{pm/K}~\cite{janz2023}.

\subsection{Unknown parameters and priors}

We add process unknowns in stages. Stage~B adds the etch depth $d$ of a partially etched rib, with the rib angle
taken equal to $\theta$. Stage~C gives the rib its own angle $\theta_r$, as when the partial etch is a separate
step. Stage~D adds the silicon and oxide indices $\nsi$ and $\nox$, and stage~E the temperature $T$, which enters
through the thermo-optic coefficients of silicon~\cite{frey2006} and silica~\cite{leviton2008}. Unless stated,
results are for stage~D. ``Realistic'' priors are $\sigma_d=\qty{2}{nm}$, $\sigma_{\nsi}=10^{-4}$,
$\sigma_{\nox}=10^{-3}$ and $\sigma_T=\qty{1}{K}$. These are assumptions; the silicon value is close to the
\num{8.67e-5} absolute index uncertainty measured for bulk silicon at \qty{1.5}{\micro m} and
\qty{295}{K}~\cite{frey2006}, while device-layer doping, stress and deposited cladding oxides are not covered by any
source we could verify.

\subsection{Why a strip cannot see its sidewall angle}

Tilting a wall by $\theta$ about its bottom edge at height $y_\mathrm{pivot}$ moves the wall inward by
$\delta x(y)=(y-y_\mathrm{pivot})\tan\theta$. To first order, the change of $\neff$ caused by moving a dielectric
boundary is a weighted integral of the displacement over the boundary~\cite{johnson2002},
\begin{equation}
\delta\neff\propto-\int_{\mathrm{wall}} W(y)\,\delta x(y)\,dy,
\label{eq:boundary}
\end{equation}
with $W(y)=\Delta\varepsilon\,|\mathbf{E}_\parallel|^2+\Delta\varepsilon\,(\varepsilon_1/\varepsilon_2)\,|E_\perp|^2$
evaluated on the core side, $\Delta\varepsilon=\varepsilon_1-\varepsilon_2$, and $\varepsilon_1$, $\varepsilon_2$
the core and cladding permittivities. A uniform width increase $\delta w$ moves both walls outward by $\delta w/2$.
Comparing the two perturbations for both walls of a symmetric waveguide gives
\begin{equation}
\frac{\partial\neff/\partial\theta}{\partial\neff/\partial w}=-h\,R\,\tan(\qty{1}{\degree})\ \text{per degree},
\label{eq:R}
\end{equation}
\begin{equation}
R=\frac{y_c-y_\mathrm{pivot}}{h/2},
\label{eq:Rdef}
\end{equation}
where $h$ is the height of the tilted wall and $y_c=\int yW\,dy/\int W\,dy$ is the centroid of the wall weight. We
write~\eqref{eq:R} in the form of a central difference over $\pm\qty{1}{\degree}$, which is how the derivatives are
computed; the exact derivative has $\pi/180$ in place of $\tan(\qty{1}{\degree})$. A tilt
therefore acts exactly like a width change of $-2(y_c-y_\mathrm{pivot})\tan\theta$.

For a strip with oxide both below and above, the structure and its mode are mirror-symmetric about mid-height, so
$y_c=t/2$ and $R=1$ exactly. A \qty{1}{\degree} tilt is then indistinguishable from a width change of
$t\tan(\qty{1}{\degree})=\qty{3.8401}{nm}$ for $t=\qty{220}{nm}$. The relation holds at every wavelength, so it
holds for $\ngr$ as well as for $\neff$, for TE and TM, and for every strip width. The angle column of $\mathbf{J}$ is
then a fixed multiple of the width column for every strip observable, $\mathbf{F}$ is singular, and no set of strips
determines $\theta$. The rotation of the wall about mid-height that remains after removing the uniform shift is odd
under the mirror operation and enters only at second order.

A rib breaks the mirror symmetry. Its wall is the upper part of the cross-section, its pivot is the slab top, and the
slab pulls the mode down toward the wall's lower corner. Then $y_c<y_\mathrm{pivot}+h/2$ and $R<1$, and the angle
column of a rib is no longer parallel to its width column. Whether this new direction survives in a joint fit depends
on how many unknowns the rib brings, which we count in Section~\ref{sec:structures}.
